\section{Dados e Metodologia}

Foram combinados os ficheiros de treino (362 perguntas) e teste (91), sem ajustar os modelos, num total de 453 perguntas de 23 temas: 230 HUM e 223 LLM. A coluna \textit{Método} indica a origem da pergunta. O carregador corrige sequências de escape do Excel, como \texttt{\_x0080\_}, e datas convertidas em cadeias; acrescenta também identificadores e a indicação do ficheiro de origem.

AMALIA (\texttt{teex/amalia:q4\_K\_M}) e Llama 3.1 8B (\texttt{llama3.1:8b}) foram consultados localmente através do Ollama. A geração usa \textit{temperature}=0, semente 42 e limite de 120 tokens; o prompt pede respostas curtas em português europeu. Um segundo modelo, Qwen 2.5 7B, atua como juiz e classifica a resposta face à pergunta e à referência. O juiz é instruído a responder exclusivamente em JSON, no formato \texttt{\{"veredito": "CORRETO"\}}; o \textit{parser} valida a estrutura, as chaves e os valores permitidos, rejeitando respostas ambíguas ou com texto adicional (e reconhece tokens exatos isolados para compatibilidade com resultados anteriores). As chamadas HTTP ao Ollama incluem até três tentativas e \textit{exponential backoff} para lidar com \textit{timeouts}, falhas de ligação e erros HTTP transitórios. As fases de geração, julgamento e pontuação são separadas e usam cache JSONL para permitir retomar a execução.

A avaliação combina EM (igualdade após normalização de maiúsculas, acentos e pontuação), Contains (referência normalizada contida na previsão), ROUGE-L F1 por tokens \cite{lin_rouge_2004}, BERTScore F1 com \texttt{bert-base-multilingual-cased} \cite{zhang_bertscore_2020} e o veredito binário do juiz. EM, Contains e ROUGE-L medem correspondência lexical; BERTScore estima semelhança contextual e seleciona automaticamente CUDA via PyTorch quando há GPU compatível, recorrendo à CPU caso contrário. O juiz é uma avaliação automática, não uma anotação humana. Os prompts estão em \texttt{qa\_pipeline/prompts.py}.
